(1) 圆形限制性三体问题中，两主星总质量、相距及共转角频率均归一化为1，质量分别为 \(1-\mu,\mu\)，其中 \(0<\mu\le 1/2\)。共转坐标下主星位于 \((-\mu,0)\)、\((1-\mu,0)\)。建立零质量粒子的共转方程及有效势 \(U(x,y)\)，严格证明 \(\mathrm L_4,\mathrm L_5\) 的位置。
(2) 推导粒子的 Jacobi 积分 \(C\)，并计算 \(\mathrm L_4\) 处势函数的全部二阶偏导数；写出围绕平衡点的完整线性化方程，指出为什么不能仅用势的 Hessian 正定性直接替代谱稳定性判断。
(3) 计算线性特征多项式及两支频率平方，求平面 \(\mathrm L_4,\mathrm L_5\) 谱稳定所需的严格质量比阈值 \(\mu_{\mathrm R}\)；讨论等号处重根的边界意义。
(4) 求 \(\mu\to0^+\) 的快、慢振动频率渐近式，讨论 \(\mu>\mu_{\mathrm R}\) 时特征根的实部；利用 Jacobi 积分的微分恒等式核验科里奥利项不改变共转积分。
